% HISTORICAL INTEGRATION SNIPPET: promoted on 29 August 2026 after the
% exact-byte production Rethlas pass for the P06.1 statement and proof.

% Replace the problem environment at \Cref{prob:frontier-pluriharmonic} by:
\begin{problem}\label{prob:frontier-pluriharmonic}
Let \(K\) be complete and nontrivially valued and let \(X\) be a good,
separated, strict, boundaryless \(K\)-analytic space.  Is every classically
pluriharmonic function \(\mathbf R\)-PL near the non-quasi-smooth points of
analytic dimension at least two?  If a continuous \(h\) satisfies
\(\dprimeddprime[h]=0\), is \(h\) classically pluriharmonic, or harmonic under
the hypotheses where a classical harmonic notion is available?
\end{problem}

% Insert immediately after the problem environment, before the paragraph
% beginning ``Weakly smooth...''.
The coefficient convention in the first question is sharp.  On the strict
semistable annulus, \(\sqrt{2}\log|T|\) is classically pluriharmonic and
\(\mathbf R\)-PL but is neither \(\mathbf Q\)-PL nor integral PL; see
\Cref{prop:p061-integral-pl-counterexample}.

The quasi-smooth scope is now settled.  By
\Cref{thm:p061-quasismooth-ph-rpl}, every classically pluriharmonic function
on a good strict quasi-smooth analytic space is \(\mathbf R\)-PL and GJR
harmonic.  The proof first passes, after a finite ground-field extension, to
Temkin's rigid \(\acute{e}\)tale cover by an affinoid with a strictly
semistable formal model.  It then descends \(\mathbf R\)-PL regularity through
the full rigid \(\acute{e}\)tale cover by the de Jong--van der Put sheaf
criterion; finite-flat descent is used only for the finite
\(\acute{e}\)tale affinoid equalizers where its hypotheses hold.  The theorem
requires no global quasicompactness, separatedness, or boundarylessness.

The published curve comparison also settles every strict analytic curve,
including singular and nonreduced curves
\cite[\S4.1.2 and Proposition~3.15]{GR25Classical}.  Hence, inside the scope
of \Cref{prob:frontier-pluriharmonic}, P06.1 remains open only at
non-quasi-smooth points in analytic dimension at least two.  The exact frozen statement and
source-aware proof passed production Rethlas with a
\texttt{gpt-5.6-sol}/\texttt{xhigh} final verdict \texttt{correct}, zero
critical errors, and zero gaps.  The authoritative identifiers are:
\begin{center}
\small
\begin{tabular}{@{}ll@{}}
controlling run & \nolinkurl{20260829T043457Z_0b1b4d161e3b}\\
controlling item & \nolinkurl{pi_7cd35ebd843cf80a7cd696bb}
\end{tabular}
\end{center}

% Replace the final summary sentence of the subsection by:
Thus P06.1 is settled on the entire quasi-smooth locus and on every strict
analytic curve, but remains open at non-quasi-smooth points in dimension at
least two.  The real-PL zero-current subproblem, the proper full-rank moment-factor
subproblem, the published smooth-proper-curve zero-current subproblem, and
the smooth-proper-curve ordinary psh regularization subproblem are settled in
their displayed scopes.  The arbitrary continuous zero-current question
remains open outside those scopes.
